% !PRINTED 407
\weberpage{407}
Dass auch $\delta = -12$ ausgeschlossen werden muss, ergiebt sich
aus der Gleichung 2., die für $\delta = -12$, $\gamma = -49$ die Congruenz
\[
\beta^2 \equiv 6 \pmod{7}
\]
zur Folge hat, die aber nicht lösbar ist.

Es bleibt also nur noch übrig:
\[
\delta = -162,\quad \gamma = 41.
\]
Aus 3. folgt für $\beta$ die Congruenz:
\[
17\,\beta \equiv 93 \pmod{145}
\]
oder
\[
\beta \equiv -1 \pmod{5},\quad 2\,\beta \equiv -1 \pmod{29},
\]
also
\[
\beta \equiv 14 \pmod{145}.
\]

Nimmt man $\beta = 14$ an, so folgt $\alpha = -8$, und alle Gleichungen 1. bis 5. sind befriedigt. Es hat sich also damit die
Zerlegung ergeben, die nachträglich leicht zu verificiren ist.
\[
\begin{aligned}
x^{10} - 36\,x^8 + 528\,x^6 - 3897\,x^4 + 14354\,x^2 - 21025 = \\
(x^5 - 8\,x^4 + 14\,x^3 + 41\,x^2 - 162\,x + 145)\\
(x^5 + 8\,x^4 + 14\,x^3 - 41\,x^2 - 162\,x - 145).
\end{aligned}
\]
